\section{Lessons Learned}

\paragraph{Taking a Stand}


When I was a young family man, I tried to be active in my community. One of those activities was to be a volunteer facilitator at a low cost group counseling centre. We underwent a Rogerian style training in being non-judgmental, reflective, and confidential. For some reason I didn't and still don't understand, two older women took a dislike to me. Furthermore, they carried it outside the training session, engaging in gossip, slander, and detraction.

When the psychologist-trainer learned of this, he was upset and reprimanded the women at a subsequent training session. In Rogerian fashion, a discussion ensued. To my surprise, the group became divided, some siding with me, and others with the hags.

I learned from that event, that there will always be those for you and those against you. From that moment on, I have never been afraid to take a stand.

\paragraph{Cash Value and the Gentleman}


Sometimes when I tell someone about my blogs, the large number of visitors, and the nearly million words of original text and translations, the response is, ``Are you making any money from it?'' Sheepishly, I have to admit the no, and I rue that a topic of conversation has been lost. I get emails about how to make money at blogs. I suppose by maintaining dozens of blogs and getting a few bucks from each, it can be done.

Most blog writers solicit tips like a waitress in a cocktail bar. The amounts are small, a five or a ten, perhaps a twenty. To accept such a tip is actually a cause for shame and embarrassment, the assumption being that such a paltry sum is important to me. The posts are freely offered although their actual value is much more than \$10.

The tradition is that a Hermetist does not make a living from the Work. Rather he would have been a horse trader, a street performer, or some such trade. Aristotle said the highest life is that of the philosopher who understands the nature of man, God, and the universe. However, first he must be a gentleman, ``noble and good'', with a certain degree of financial independence.

Thus, this is not a ``blue collar'' blog with crude language, irrational views, the promotion of immoral lifestyles, and ugly presentation as a few of you persist in believing have any value for a gentleman.

Rather than money, I would like some among you to pray for my health, the good birth of a grandson, and a successful outcome for my son's upcoming surgery.

\paragraph{A Prophet without Honor}


It may or may not surprise you that no one, literally, who knows me personally reads this blog; this includes family and close friends. They will not even ``like'' the Gornahoor facebook page. I am the ``crazy uncle'' that the children like to be with, with outlandish yet ignorable views; my sister calls me a shaman.

There are two possible conclusions that readers can draw from this. The first is that those who know me best think I am a ``kook'', so you should certainly not pay attention.

The other conclusion is that if you try to step up and out in your life, your friends and family will be the first to try to hold you back.

\paragraph{Runes in Salemi}


A half dozen years ago, I went to Salemi, Sicily, to show my son our family roots and to look for property to buy. The city is built on a hillside overlooking olive groves, vineyards, and farmland. At the top of the highest peak, there is a fort with a view for many miles around. Originally built by the Arabs, it was eventually taken over by the Normans. I was hoping to find evidence of Norse culture there, perhaps even Rune symbols, as well as the much larger fort in Palermo.

I did find a good property, three stories, with kitchens on the 2\textsuperscript{nd} and 3\textsuperscript{rd} floors, and enough room to sleep 8 or even a dozen. I was hoping to found a centre there for those interested in the topics discussed on Gornahoor. With a European centre, it would have been possible to work in groups where I could have taught actual spiritual exercises. Unfortunately, the economy turned, depleting my investments, and an inheritance was much less than expected.

Of course, I found no rune inscriptions, since the Normans had long since abandoned those ways. However, it was in Palermo that I was struck with a deep understanding. These Normans had moved into France, conquered England, Sicily, and from there, went on to take over leadership the Holy Roman Empire. I wondered about the inner state of these knights who had come so far and accomplished so much. In the tour of the castle in Palermo, we eventually made it to the Palatine Chapel in it and there I found the answer. It was certainly beautiful and I could sense the knights of the past attending mass in it.

Since then I have no patience with those who deny the spirituality of those knights in favour of some remote, imaginary paganism. Particularly, since they are usually crude, vulgar, and ungentlemanly about it.

\paragraph{The Limits of your Imagination}


At a now defunct minicomputer manufacturer, we had developed a networking system; at that time, there were no turnkey solutions as there are now, so it was a unique accomplishment. At one meeting, we had to listen to a presentation of an earnest young women who was tasked with marketing the networking product.

Unsure of what to say, she finally blurted out with unfeigned enthusiasm, ``And the possibilities for this product are limited only by your own imagination.''

Of course we all mocked her. What that meant is that she could not articulate any creative or practical uses for the network, so she was counting on each audience member to think one up for her.

That is when I learned the lesson of specificity. It is difficult to provide specifics, so most people just fake it, counting on the audience to ``know'' what they mean. Thus we hear about a vague, unspecified ``Indo-European tradition'', not unlike Chomsky's universal grammar; of course, you cannot specify it. Curiously, we have just written about this, but some readers don't know how to learn a lesson.

\paragraph{The Image of an Image}


When the kids were young, I had to take them to Disneyworld every now and then. There was a last time, after which I got too disgusted to return. We were on an ``adventure'' ride, and a mechanical deer would randomly pop up out of the ground. The passengers would ooh and ah, pointing to their children.

When you see through the contrivance, it has no effect. I thought, why not go into the woods and look for a real deer? If nature is the image of the divine, as the holy ones teach us, then what Disney provides is merely an image of an image. Yet, on a larger scale, most of our lives today consists of images of an image, from movies, to TV, to the popularity of social media. Nothing good can come of it.

\paragraph{Actuarial Science}


In high school, I was one of the top two dozen scorers in a Massachusetts statewide math exam, so I was invited to meet with its sponsor, the Society of Actuaries. There were follow ups and I seriously considered a career in the field. By reducing everything to quantity, it manages to provide concrete and successful results. I learned some of its techniques, and still think it is applicable for many tasks of the producing class.

For example, as you all know, the insurance rates for teenage boy drivers are much higher than for anyone else. That is a simple matter of actuarial science which I cannot explain here. Some people may consider that unfair. However, when an insurance company publishes its rates, it does not have to provide copious and apologetic footnotes along the lines that ``not every teenage boy is a bad driver'', or ``not every bad driver is a teenage boy'', etc.

The simple mathematical fact is that without that policy, the insurance companies would simply go out of business. Contemporary western societies today are ``going out of business'', but it is happening too slowly for most to notice.

\paragraph{A Continuous Tradition}


A few months ago after the Venner controversy, we lost some readership. I read his blog and he indicated his disappointment that France did not have a religion like Hinduism to preserve its racial and ethnic identity. As an historian, he had to know that France had no problem preserving its identity prior to the Revolution. However, as an ideologue, he could not admit this. Ideology is the image of an image in the intellectual realm, i.e., the image of a thought which is the image of the spirit. Nothing good can come of that, but the world is full of ideologues.

\paragraph{Camel in the Desert}


There is a recent comment from someone who is expecting the ``real'' Indo-European tradition to show up, presumably just for him. To boot, it is in a post that warns against seeking any perfection in this world.

Does a man dying of thirst in the desert argue about which camel will take him to the oasis? Does he languish there waiting for the perfect camel to come along? There is a lesson there that few will heed.

\paragraph{Appeals to the Intellect}


Let us note again the three forces in the psyche: sex, anger, and intellect. Since these are fundamental, the vices of the first two are, as Augustine pointed out, insatiable. That is why nearly all media are geared to stimulating sexual desire or anger. The former is too common to provide examples. The latter is the basis for all the political blogs and talk shows in our time. They ``work'' only if they get you worked up about some issue, person, place, or thing. The lesson is to always be on the lookout for that.

Beyond that, there is the appeal to the emotions in general. Every show emphasizes it. News stories depend on them for the ``human angle''. Along with that, there are always stories about dogs, elephants, or other animals displaying ``human-like'' emotions. These are very popular.

However, the truth is just the opposite. The real lesson is that, for the most part, human emotions are more animal-like and that explains the rapport. That is not merely metaphysical, but also scientific as described in \textbf{Carl Sagan}'s book, \emph{The Dragons of Eden}. Now I am not at all opposed to building rapport with animals, but only to the idea that it somehow represents what is higher in man.

Appeals to the Intellect are much rarer, and that is necessarily so. Few can attain the equanimity of Aristotle's philosopher. The philosopher, meant in the ancient sense as a way of life, would want to see into the nature of things and passions can only work against this. Even writings that aspire to a higher level are still stuck on an ``us-vs-them'' attitude, which can never transcend the thymos.

The philosopher is the lover of wisdom, but it is still not the state of being wise. When Jesus said we must love our enemy, it means we need to transcend all points of view and see things as they are.


\flrightit{Posted on 2013-10-17 by Cologero}

\begin{center}* * *\end{center}

\begin{footnotesize}
\begin{sffamily}
\texttt{X on 2013-10-18 at 01:22 said:}

Many valuable perspectives.

\hfill

\texttt{Scardanelli on 2013-10-18 at 06:24 said:}

``The shaman'' doesn't sound half bad... At age 32, I am known as ``the grandpa'' by my wife and kids, who unfortunately are thoroughly under the sway of this world. Unfortunately it seems that the closer a person is to you the less they'll listen to you.

I am grateful for your work, and I will pray for you and your family.

\hfill

\texttt{Michel on 2013-10-18 at 07:35 said:}

``I did find a good property, three stories, with kitchens on the 2nd and 3rd floors, and enough room to sleep 8 or even a dozen. I was hoping to found a center there for those interested in the topics discussed on Gornahoor. With a European center, it would have been possible to work in groups where I could have taught actual spiritual exercises. Unfortunately, the economy turned, depleting my investments, and an inheritance was much less than expected.''

Yeah, that would have been a great idea, those darn Bavarians with their secret societies! Ruining everything for everyone, tsk!

I do hope that you get well soon, and that those with the same tinge of madness you have eventually see that you are the grand court jester, the keeper of deep truths, the heyoka who shivers when it is sunny, dresses lightly during winter, in being yourself, you are the real them. It's interesting, the last four posts of yours have been strangely consistent with some of the things that are happening in my life, not only in terms of thoughts, but actual situations and points where I have to make `strange' decisions, this very post is rather interesting, because I've just had a discussion with my father on certain things that are to him, `kooked-up' before reading this. Even more interesting is a talk I was having with an Italian friend of mine on Wednesday evening, he is an initiate in a Sufi Order, the Qadirriya Chistiyya, I know that this blog aims at pointing towards Western Tradition, so I mean no disrespect at all by speaking on this, but he said something similar to what you said when you said this:

... 

''

Does a man dying of thirst in the desert argue about which camel will take him to the oasis? Does he languish there waiting for the perfect camel to come along? There is a lesson there that few will heed.

''

... 

I was telling him of your efforts when he said that on the one hand, it is good what you are doing (not that you need anyone's approval), but he also fears that you may hinder others from turning to Traditions that are `Alive' inspite of their being Eastern, such as Islam. I asked him about the Law of Affinity, in that it is the Truth, manifested physically as an Order that selects its members via this Law, in that those with similar `vibrations' to it would be attracted to it naturally, and that they truly would be qualified to join it, whether it was hidden or in the open. I said hidden or in the open because I read here, and have also given thought on the possibility that certain Western Traditions did go `underground' after the Templars were chased away, and that their `Operation' is of a different stride or manner, in that it, more than the Eastern Orders, due to the current degenerate Western condition, works solely along inner lines, sort of like an `inner Theophanic collective consciousness'. We also spoke about the different Orthodox writings of the Greek Church and the various monasteries that still exist today, about Freemasonry and its `Sufi Origins' (note my use of quotation marks on that) about Hermeticism and also about Roman Catholicism. We agreed that while it would be possible that there may exist in the West an Order that maintains an unbroken initiatic chain, he said that ``God has given you a river right before you, but you would rather go to the desert in search of an oasis to quench your thirst''. He also says that the other Traditions are like the stars, but Islam is like the sun, in that being the final Divine Revelation, it is most suitable for the current world conditions and it also synthesizes all the others, he further said that all the avatars and prophets were all pre-figurations of The Prophet, and that he encompasses all of them, being the Initial Ray that proceeds from God, such that in Essence, he is God, the ray essentially is the sun, from whence it derives its very nature, it is the sun `extended'. I'm not proselytizing, it's all about going beyond the form, riding the camel till it brings you to the oasis, the destination is the oasis, not the camel's back. Anyways, I pray that all your life's affairs turn out the way they were meant to.

\hfill

\texttt{Tosti on 2013-10-18 at 08:21 said:}

Your thoughts are always appreciated and, as you know, they do have effect.

\hfill

\texttt{Michael on 2013-10-18 at 10:00 said:}

A shaman? That would make you very popular among fans of ``Eat, Pray, Love'' but I think you are closer to what is described in letter IX.

You can count on my prayers for you and your family.

\hfill

\texttt{visionsofglory14 on 2013-10-18 at 11:28 said:}

I know you personally, and I read your blog.

\hfill

\texttt{Thorsten on 2013-10-18 at 12:04 said:}

``Since then I have no patience with those who deny the spirituality of those knights in favour of some remote, imaginary paganism. Particularly, since they are usually crude, vulgar, and ungentlemanly about it.'' I've made it clear that I am not denying the spirituality of those knights or of any of the other knights of that era. On the utter contrary... in order to defend their their spirituality I think a Transformation would have to make it clear that this is the western Tradition, in other words the Indo-European Tradition. To that end would it not be necessary for the symbolism, that is to say the exoteric shape of Christianity, which has evidently lead so many astray to be adjusted? It has been repeated to me that Guenon and all other traditionalists authorities make it clear that a religion must undergo a spontaneous transformation (correct me if I'm wrong) by revelation. I understand this and am awaiting it just like anyone else who is convinced of this truth. I've simply indicated what I believe will be one or two of the exoteric characteristics of such a transformed faith. Why this modest assertion has been so misunderstood eludes me.

``That is when I learned the lesson of specificity. It is difficult to provide specifics, so most people just fake it, counting on the audience to ``know'' what they mean. Thus we hear about a vague, unspecified ``Indo-European tradition'', not unlike Chomsky's universal grammar; of course, you cannot specify it. Curiously, we have just written about this, but some readers don't know how to learn a lesson.''

I'm not aware of any one who has done research to determine the underlying metaphysics represented in Indo-European myth and cosmologies but it should be possible in theory. I assumed most here had at least heard of Indo-European studies (\url{http://en.wikipedia.org/wiki/Indo-European\_studies}) which is a natural companion science to studying Traditionalism in the context of European culture. Evola was evidently pretty well read in the literature in this field during his lifetime at least as far as the writings of Georges Dumézil are concerned. I don't see why suggesting the relevance of this discipline is out of place, let alone indicative of any attempts to obfuscate. In the interest of the spirit of free inquiry which I recall being written of positively here not too long ago, I would expect I.E. studies to be of interest in a place where the ``hyperborean tradition'' has been referenced on so many occasions.

``The philosopher is the lover of wisdom, but it is still not the state of being wise. When Jesus said we must love our enemy, it means we need to transcend all points of view and see things as they are.'' I've indicated the possibility of a coming religion which may arise that transcends the viewpoint of one historical manifestation of the divine, while gathering into itself other manifestations of the divine. In sum a religion in which Christ is one avatar among many, and in which it is no longer possible to mistake his charity and sacrifice as some form of ``Bolshevism of the ancient world'', as it was nothing of the sort. Would this not be a solar view point transcending any illusory false dichotomies? I thought this suggestion would be received as ``irenic'', rather than ``eristic'' approach.

In any event, Cologero the best of health to you and your son, and the best of luck on the birth of your grandson. Your work is appreciated very much and has made a deep impact on me, and will continue to do so.

P.S. I look forward to discussing the symbolism of the New Testament when I'm done reading the Jung text.

Sincerely,

Thorsten

\hfill

\texttt{Ash on 2013-10-18 at 19:32 said:}

Thank you for this post. Understandably, we want to in many ways minimize our ``personal equations'' when undertaking this study, but I certainly benefit from hearing about the experiences that led you and others here. My own blog has fewer readers than yours, I'm sure, in many ways due to the somewhat chaotic posting schedule. I am always impressed that you are able to write as voluminously as you do.

Speaking for myself, the reason I abandoned an economics/political science combined major for my current focus in food systems and agricultural economics is precisely the observations you yourself reported. I grew tired of the ``fluff'' that came with them, even in philosophy. I can work with something solid, with results. I can earn a living that, Providence willing, will be based on the Traditional principles manifesting in what I do. When I discuss philosophy or politics with friends in those disciplines, I'm doubtless the ``kook'' in the situation. But then, the holy fool has a very unique freedom.

I was just recently criticized by a friend for being too ``aggressive'' in my argumentation, picking at arguments in ways that come across as arguing for the sake of argument. Strangely enough, it came up in the conversation that the friend considered moral philosophy to be, in the end, just a matter of interest which in the real world paled in comparison to other factors such as moral intuitions. Your last words in this post are undoubtedly worth pondering.

I will keep you and your family, born and unborn, in my own prayers.

@Thorsten: The future religion you describe sounds like it has some similarities to the Baha'i heresy, which also holds that Christ, Muhammad, Krishna, and others were ``equal avatars.'' Could you explain what you hold Christ to be? Is he unique in being the Logos and as ``Messiah''? Are other avatars incarnations of other aspects of the Divinity?

\hfill

\texttt{Cologero on 2013-10-24 at 00:35 said:}

I was referring to ``Carl Rogers'', who was an influential psychologist back in the dat.

\hfill

\end{sffamily}
\end{footnotesize}

